\documentclass[9pt]{extarticle}
\usepackage[a4paper,margin=1.5cm]{geometry}
\usepackage[T1]{fontenc}
\usepackage[utf8]{inputenc}
\usepackage{times}
\usepackage{microtype}
\usepackage{multicol}
\usepackage{enumitem}
\usepackage{tabularx}
\usepackage{booktabs}
\usepackage[hidelinks]{hyperref}
\usepackage{titlesec}
\usepackage{parskip}
\usepackage{amsmath}
\usepackage[round,authoryear]{natbib}

% Version recortada de dos paginas (entrega 2026-09-18). Responde a la
% retroalimentacion sobre la primera version (PFCII___ENG___Kiara.pdf).
% Bibliografia: ../refs.bib, generado con scripts/build_refs.py. No editar aqui.

\titleformat{\section}{\normalfont\bfseries\large}{\thesection.}{0.5em}{}
\titlespacing*{\section}{0pt}{5pt}{2pt}
\setlist{nosep,leftmargin=*,labelsep=0.5em}
\setlength{\parskip}{3pt}
\pagestyle{empty}

\begin{document}

\begin{center}
  {\Large\bfseries MetalSynth-Pelvis: Conditioned synthesis of osteosynthesis implants\\[2pt]
  and local metal artifacts in pelvic CT scan volumes using Multi-Window Latent Diffusion}\\[3pt]
  {\normalsize Kiara Alexandra Balc\'azar Santa Cruz --- UTEC \quad $\cdot$ \quad Advisor: V\'ictor Flores Benites}
\end{center}
\vspace{1pt}\hrule\vspace{4pt}

\section{Problem Statement}

Segmentation networks for pelvic computed tomography (CT) learn from annotated volumes, and few annotated volumes contain osteosynthesis implants. In CTPelvic1K \citep{liu2021ctpelvic1k}, the metal-bearing subset (CLINIC-metal) carries few bone annotations. Metal causes beam hardening, streaking and photon starvation, and these artifacts extend beyond the implant into the bone around it. Synthetic implants could supply such cases. This thesis builds and evaluates the synthesis method. It does not measure the effect on segmentation.

Three lines of work address parts of the problem. Diffusion models for lesion synthesis, such as those of \citet{chen2024tumorsynthesis}, \citet{zhang2025diffboost} and \citet{ramzan2026claim}, condition on a mask and contain no mechanism or training signal for intensity changes outside the object, whereas metal changes intensity outside itself. Physics-based simulation with CatSim \citep{deman2007catsim} within XCIST \citep{wu2022xcist}, as used in the hybrid protocol of \citet{peters2025hybrid}, places metal at random positions for training data and by expert choice for the benchmark. It does not place implants under anatomical constraints. Surgical planning pipelines \citep{liu2025pipeline} compute one trajectory, while clinical placements follow a distribution: on a four-grade cortical-breach scale \citep{smith2006iliosacral}, \citet{zwingmann2009navigated} report grade~0 in 69\% of navigated and 40\% of conventional iliosacral screws.

\textbf{Gap.} No published method combines (C1) rigid parametric implant geometry, (C2) a placement sampler that matches measured malposition distributions, and (C3) multi-window HU encoding, taken from metal artifact reduction \citep{wang2025adaptiveweighting}, applied to synthesis, with a generation band $B_{\delta}$ (about 12~mm) beyond the implant mask so that streaks and beam hardening appear outside the metal.

\section{Research Question and Hypothesis}

\textbf{Question.} Compared with copy-paste insertion and with the hybrid physics protocol of \citet{peters2025hybrid}, does a placement sampler constrained by bone density, cortical clearance and corridor viability, combined with a multi-window latent diffusion model conditioned on implant geometry and $B_{\delta}$, produce pelvic screws and local artifacts with higher surgical and physical coherence?

\textbf{Hypothesis.} Sampled placements lie closer (Wasserstein-1) to the two S1 breach-grade distributions of \citet{zwingmann2009navigated} than unconstrained placements. Synthesized appearance differs less from real artifacts than copy-paste insertion, and its streak amplitude matches that of the physics protocol.

\section{Objectives}

\begin{enumerate}
  \item \textbf{Representation gate (Go/No-Go).} Round trip from HU to three windows, through the Stable Diffusion 1.5 autoencoder \citep{rombach2022latentdiffusion}, and back to HU, with bone mean absolute error (MAE) below 25~HU on 34 held-out patients. Two autoencoder variants (pretrained, and decoder adapted with encoder frozen) and three encodings are tested. The rule was fixed before any run. The 25~HU level follows the error of current methods on the same scale (RMSE of 12.3 to 20.2~HU in \citet{karageorgos2024ddpm} and 12.74~HU in \citet{yun2026simulationdriven}). If no combination passes, Objective~3 is not pursued and the No-Go is the result.
  \item \textbf{Placement sampler.} 3D pose of a parametric screw (6.3 to 8~mm) in the CT reference frame of \citet{kaiser2014dysmorphism}, constrained by bone density read from CT numbers \citep{arand2019pelvicring}, 5~mm cortical clearance and corridor viability \citep{mclaren2021corridor}, measured on TotalSegmentator masks \citep{wasserthal2023}. Target: the two ordinal distributions of \citet{zwingmann2009navigated} at S1. S2 is reported without a clinical target.
  \item \textbf{Synthesizer.} 2.5D ControlNet \citep{zhang2023controlnet} on Stable Diffusion 1.5, with frozen encoder and the generation band $B_{\delta}$.
  \item \textbf{Evaluation.} Surgical Admissibility of Placement (SAP), the only metric this work introduces, compared by Wasserstein-1 distance. Appearance uses bone integrity, metal integrity and streak amplitude under their published names \citep{peters2025hybrid}. RMSE and SSIM outside $B_{\delta}$ check that no other region changes.
\end{enumerate}

\section{Expected Results}

\begin{table}[h]
\centering
\small
\begin{tabularx}{\linewidth}{@{}p{2.9cm} p{5.2cm} X@{}}
\toprule
\textbf{Goal} & \textbf{Metric} & \textbf{Expected outcome} \\
\midrule
Representation (Obj.~1) & Bone MAE in HU after the round trip, mean per patient on 34 test patients & Below 25~HU for at least one pre-specified combination, or a reported No-Go. \\
Placement (Obj.~2 and 4) & SAP: four-grade cortical breach \citep{smith2006iliosacral}, density-zone fraction, corridor viability \citep{mclaren2021corridor} & Lower Wasserstein-1 distance to the two S1 distributions of \citet{zwingmann2009navigated} than unconstrained placement. S2 reported without a clinical target. \\
Appearance (Obj.~3 and 4) & Bone integrity, metal integrity, streak amplitude \citep{peters2025hybrid} & Lower discrepancy than copy-paste insertion. Streak amplitude comparable to the physics protocol. \\
Preservation & RMSE and SSIM against the original CT, outside $B_{\delta}$ & No change outside $B_{\delta}$. \\
\bottomrule
\end{tabularx}
\end{table}

\section{Data and Metal Extraction}

\textbf{Data.} 178 local CTPelvic1K volumes (103 in dataset6 and 75 in dataset7) from 168 patients, split by patient (126 train, 8 validation, 34 test) so that no patient enters both training and evaluation. Challenge data \citep{haneda2025aapm} are not used for training.

\textbf{Handling false positives in metal extraction.} A threshold on HU returns voxels that are not implant: cortical bone, objects outside the body, and bright streaks next to the metal. Thresholded masks over-cover the metal \citep{xie2024implantsegmentation}. Four measures address this.
(i)~\emph{Screening ranks, review decides.} A 2500~HU threshold selects volumes for review: it flagged 113 of 178 and missed none of the volumes with metal or foreign objects in a 3D review. At 1500~HU it flagged 177, since cortical bone exceeds that value. Each flagged volume is labelled by 3D and multiplanar review, which separates pelvic osteosynthesis, other implants and objects outside the body. A flag without a confirmed object does not enter the metal cohort.
(ii)~\emph{No extracted geometry enters synthesis.} Of 57 elongated objects extracted by HU threshold, 9 were fragmented and 49 measured below 6.0~mm in diameter, under the published screw calibre. Implants are therefore parametric cylinders, and extraction errors cannot reach the generator.
(iii)~\emph{Two rules bound each real metal mask.} Where a real mask is needed (space occupied by an existing screw), a local half-maximum threshold per object and a fixed 2500~HU threshold are both applied, and corridor viability fell from 53.8\% to 40.4\% of 52 patients under both. Metal integrity uses the threshold per region of interest defined by \citet{peters2025hybrid}.
(iv)~\emph{Artifacts in anatomical masks.} Components below 0.1\% of their structure are removed and each case passes quality control. A landmark region counts as contaminated when its dark-streak or metal fraction exceeds the maximum found in 69 metal-free patients. Of 65 patients with orthopaedic metal, the reference frame was computable in 48 and free of contamination in 29. Contaminated cases are reported apart and not pooled.

\section{Computational Resources and Budget}

\textbf{Resources.} Khipu, the UTEC cluster: two nodes with one NVIDIA RTX A6000 (48~GB) each, one A100 node with MIG partitions, and a CPU node for steps without GPU. The thesis account allows 3 running jobs, 32 CPU cores, 98~GB of RAM and 24~h per job. Scripts save checkpoints and resume across jobs.

\textbf{Measured costs.} Anatomical segmentation of the local cohort: 6~h~27~min. Objective~1 decoder adaptation: 30{,}000 steps at 1.32~s per step and 24.6~GB of GPU memory. One encoding finished in 10~h~36~min, so the three encodings take about 33 GPU-hours. Evaluation takes about 65~s per volume, model and encoding.

\textbf{Pending estimate.} ControlNet training time is not yet measured, and no figure is given for it here. It will be fixed by the procedure used for Objective~1: a timed pilot of 200 steps on an RTX A6000, and the number of planned steps multiplied by the seconds per step, before the full run. If the Objective~1 gate fails, ControlNet is not trained. The physics arm runs on a reduced patient subset. \citet{wu2022xcist} report a run time of ``a few minutes on a PC with 12 CPU-cores'', so the arm is planned on the CPU node, and its time per volume will be measured in its first run.

\textbf{Scope limits.} Downstream segmentation (Dice, HD95) and ablations per constraint are out of scope, which removes their training cost. Joint prostheses are excluded: in training sinograms generated by XCIST, the metal trace ``may look distorted for large metal objects with diameters larger than 3.0~cm'' \citep{haneda2025aapm}. This work models screws.

\section{References}
\vspace{-6pt}
\scriptsize
\renewcommand{\bibsection}{}
\setlength{\bibsep}{0.6pt}
\setlength{\bibhang}{8pt}
\begin{multicols}{2}
\bibliographystyle{abbrvnat}
\bibliography{../refs}
\end{multicols}
\end{document}
